\section{Discussion}

The gate does not track better; it declines to commit. Every policy in
Table~\ref{tab:reconnection-main} below the unconstrained row accepts fewer
reactivations, and the HOTA gap is the price of that refusal. Keeping the
distinction matters for how the method should be read. HOTA scores a
trajectory; the gate scores a decision. The two agree only when every
decision the gate refuses would also have been wrong, which is not what the
numbers show: the gate refuses some correct returns along with the wrong
ones, and the recovery counts are there to say how many.

Why the product terms help at one operating point and not at the other four
is the question the ablation leaves open. A plausible account is that a
relaxed ceiling is where the gate has to adjudicate borderline returns, and
borderline returns are where motion agreement decays with gap length; at a
strict ceiling the admitted set is small and already unambiguous, so no
weighting scheme has much left to separate. We offer that as an
interpretation rather than a tested mechanism. Separating the two would
require proposals stratified by their true difficulty, and our labels
record only whether a committed decision turned out to be right.

The cost is worth stating in the units of the application rather than in
HOTA points. A system that counts people, measures dwell time, or hands
trajectories to a human reviewer keeps paying for a wrong identity long
after the frame that caused it; for such a system, trading 136 recovered
gaps for 1,397 avoided failures is a good exchange. A system judged on HOTA
alone should not make that trade, and Table~\ref{tab:reconnection-main}
says so plainly.

Appearance is the obvious complement rather than a competitor. Where a
re-identification model and usable crops exist, a descriptor answers
questions this gate cannot: whether two boxes show the same clothing, the
same gait, the same carried object. What the gate offers instead is a
version of the decision that needs none of that infrastructure, together
with a measurement of how far geometry alone gets. A future study that puts
both signals on the same proposals would separate the failures caused by
missing evidence from the failures caused by a score that weighed the
available evidence badly. Our results say the second category is real and
worth roughly a third of the failure mass; they say nothing about the first.

Two further questions follow from the scope. The first is whether the
ordering survives on sources no earlier round of the analysis has touched.
The second is whether the same construction transfers past the cap of
twelve admissions per source, which is a design parameter rather than a
measured quantity; Table~\ref{tab:caps} shows the ordering holding across
caps while the magnitude moves, and a deployment would have to choose the
cap from its own tolerance for missed returns rather than from this table.
